\begin{lemma}[Immediate extensions in the front tree]
Let $F$ be a front on $X$.  If $s\in T(F)$ but $s\notin F$, then the type
\[
  E_s=\{n\in\mathbb N\mid
  \mathsf{ProperInitialSegment}(s,s\cup\{n\})
  \text{ and }s\cup\{n\}\in T(F)\}
\]
of ordered one-point extensions of $s$ in the prefix tree is nonempty and
countable.  The proper-initial-segment condition means that the new point is
the next point above the finite prefix $s$, as in the paper's notation
$n\in\omega/s$.
\end{lemma}
\begin{proof}
Since $s\in T(F)$, \noderef{PrefixTree} gives $t\in F$ such that
$s$ is an initial segment of $t$.  Because $s\notin F$, prefix-freeness in
\noderef{Front} implies that $s\ne t$, so $s$ is a proper initial segment of
$t$ in the sense of \noderef{ProperInitialSegment}.  Let $n$ be the least
element of $t\setminus s$.  The initial-segment description of finite sets
from \noderef{InitialSegment} shows that every element of $t$ below $n
already belongs to $s$; hence
$s\cup\{n\}$ is an initial segment of $t$.  Therefore it belongs to
$T(F)$ by \noderef{PrefixTree}.  The same description from
\noderef{InitialSegment} shows that
$s$ is a proper initial segment of $s\cup\{n\}$, so $n$ witnesses membership
in $E_s$.  Finally $E_s$ is a subtype of $\mathbb N$, so it is countable by
the standard countability of $\mathbb N$ and closure of countability under
subtypes.
\end{proof}
